\section{Conclusion}

Purpose-selective execution asks for more than hiding data or deleting it: a
confidential attribute should remain useful where approved and irrelevant where
prohibited. This paper defines a bounded way to test that requirement across
data, workloads, capabilities, releases, privacy, lifecycle, and evidence.

The diagnostic evaluation demonstrates that the PSBE-Runtime pipeline executes
correctly against three commercial models via OpenRouter. The admission gate
rejects models with unstable route assignments, the position diagnostic
identifies layout-sensitive release behavior, and the confirmatory matrix
validates end-to-end execution across datasets and conditions. Release denials
are concentrated in one model on one dataset.

A reduced-scope confirmatory run with 40 pairs per dataset shows 79.7\% pass
rate for the admitted model, with dataset-specific release behavior. Security
evaluation via test-double oracle shows 75.7\% attack detection rate and
100\% evidence coverage. All metrics are computed from test-double oracles;
live execution is required for definitive claims.

Three diagnostic claims are supported by frozen evidence. The full confirmatory
claim family (H1--H10) requires the registered 200-pair run with statistical
inference. All artifacts, budgets, and claim traces are frozen.
